
% =====================================================================
\section{T5: selection requires a gate}
\label{sec:gating}
% =====================================================================

A contact graph fixes which steps are \emph{possible}---the edges---but not which
step is \emph{taken}. We show that whenever a vertex has more than one forward
option, the next step is underdetermined by the graph alone, so that selecting a
walk requires an additional datum: a gating function. The necessity of such a
datum, and the form of the walk it produces, is T5.

\subsection{Underdetermination of the next step}

\begin{lemma}[Forward branching]
\label{lem:branching}
Let $\Gph$ be a contact graph and $v\in\V$ with degree $\deg(v)\ge 2$. Then there
are at least two distinct walks of length one from $v$, and the graph alone does
not determine which is taken.
\end{lemma}

\begin{proof}
By definition $\deg(v)=\abs{\Nbr(v)}\ge2$, so there are distinct
$u_1,u_2\in\Nbr(v)$ and distinct one-step walks $(v,u_1)$, $(v,u_2)$, each
admissible since $vu_1,vu_2\in\E$. Nothing in $(\V,\E,\wt)$ privileges one over
the other: both are equally edges of equal-or-incomparable weight, and the graph
carries no ordering of incident edges. Hence the next step is not a function of
$\Gph$ alone.
\end{proof}

\begin{remark}[Most vertices branch]
\label{rem:branch-generic}
A connected contact graph with $\abs{\V}\ge3$ has at least one vertex of degree
$\ge2$ \textup{(}else it is a single edge\textup{)}; in a graph rich enough to
individuate many parts, branching is the generic condition. So
underdetermination is not an edge case but the normal situation a walk faces.
\end{remark}

\subsection{The gate as the minimal selecting datum}

\begin{theorem}[Necessity of a gate]
\label{thm:gate-necessary}
Let $\Gph$ be a contact graph in which some vertex has degree $\ge2$. Then a walk
that is determined step-by-step requires a gating function
\textup{(}\Cref{def:gate}\textup{)}: there is no rule depending on $\Gph$ alone
that selects a unique next step at every branching vertex. Conversely, any rule
that does select a unique next step at each vertex \emph{is} a gating function
that selects \textup{(}assigns positive forward weight to exactly one incident
edge\textup{)}.
\end{theorem}

\begin{proof}
\emph{(Necessity.)} At a branching vertex $v$ (degree $\ge2$), \Cref{lem:branching}
exhibits $\ge2$ admissible next steps not distinguished by $\Gph$. A rule that
nonetheless outputs a unique next step must use a datum beyond $(\V,\E,\wt)$ to
break the tie. Encode that datum as $\gate(v,\cdot)$ by setting
$\gate(v,e)>0$ for the chosen edge and $0$ for the others; this is a gating
function in the sense of \Cref{def:gate} (nonnegative, sub-normalised). Hence a
gate is required and the selecting rule is realised as one.

\emph{(Converse.)} If $\gate$ assigns, at each vertex, positive forward weight to
exactly one incident edge, then at every step the admissible continuation is
unique, so $\gate$ determines a unique walk from any start vertex. Thus selecting
rules and selecting gates coincide.
\end{proof}

\begin{definition}[Selected walk; the gated trajectory]
\label{def:gated-walk}
Given a selecting gate $\gate$ and a start vertex $v_0$, the \emph{gated
trajectory} is the unique walk $\Walk_\gate(v_0)=(v_0,v_1,\dots)$ with each
$v_{i+1}$ the $\gate$-selected neighbour of $v_i$. When $\gate$ assigns positive
weight to several edges \textup{(}a \emph{soft} gate\textup{)}, it determines
instead a distribution over walks, with the probability of a walk the product of
its step-weights; the selecting case is the deterministic limit.
\end{definition}

\subsection{The gate biases an otherwise free graph}

\begin{proposition}[Reachability is near-total; the gate fixes the order]
\label{prop:reachability}
In a connected contact graph, every vertex is reachable from every other by some
walk. Hence the graph alone does not constrain \emph{which} vertices a trajectory
visits, only that any are reachable; it is the gate that fixes the \emph{order}
in which they are visited and therefore the trajectory's identity.
\end{proposition}

\begin{proof}
Connectedness gives a path between any two vertices (\Cref{ax:finite}), so all
vertices are mutually reachable; the set of reachable vertices from any start is
all of $\V$. Thus the reachable set is invariant (everything), while the realised
walk depends entirely on the step-choices, i.e. on $\gate$ (\Cref{def:gated-walk}).
Two different gates from the same start in general yield different walks through
the same reachable set, so the gate, not the graph, individuates the trajectory.
\end{proof}

\begin{remark}[Interpretation: the selecting field, on which nothing depends]
\label{rem:field}
One may read the gate as a field that, given a situation, biases which step a
part takes next---so that what is ``thought next'' from a given vertex is set not
by the graph (which makes everything reachable, \Cref{prop:reachability}) but by
the gate. In that reading the order, not the availability, is the content of a
trajectory, and the gate is the situation-dependent selector that supplies the
order. Different situations are different gates and hence different orders over
the same reachable graph. We use only the graph statement
\textup{(}\Cref{thm:gate-necessary,prop:reachability}\textup{)}; the reading is
orientation and is not relied upon.
\end{remark}

\begin{remark}[Why a gate is cheaper than the graph]
\label{rem:gate-cheap}
A gate is a local datum---one sub-normalised distribution per vertex---whereas
specifying a walk by exhaustively comparing all global continuations would
require resolving the whole graph at each step. The gate is the minimal addition
that makes stepping decidable: it is to the trajectory what a local rule is to a
global search. This minimality is why selection is realised by gating rather than
by re-deriving the graph at every step.
\end{remark}
